\chapter{Technology Stack}
\section{Frontend Stack}

The frontend of the system is designed to provide a responsive, interactive, and user-friendly interface. It emphasizes modularity, reusability, and maintainability of components while ensuring fast load times and smooth navigation.

\subsection{Framework and Libraries}
\begin{itemize}
    \item \textbf{React.js} --- Used as the primary framework for building a component-based user interface. React enables efficient rendering through its virtual DOM and promotes reusability of UI components.
    \item \textbf{Vite} --- Adopted as the build tool and development server. Vite provides a faster development experience with hot module replacement and optimized builds.
    \item \textbf{CSS} --- Utilized for styling to achieve a clean, modern, and responsive design with utility-first classes.
    \item \textbf{Lucide Icons} --- Used for lightweight, customizable icons across the interface.
\end{itemize}

\subsection{Design Principles}
\begin{itemize}
    \item \textbf{Responsiveness} --- Ensures seamless user experiences across devices with different screen sizes.
    \item \textbf{Accessibility (aria)} --- Follows best practices to make the application usable for all users, including those with disabilities.
    \item \textbf{Performance Optimization} --- Lazy loading, code splitting, and optimized assets improve loading times.
\end{itemize}

\subsection{Testing and Quality Assurance}
\begin{itemize}
    \item \textbf{Playwright } --- Applied for integration testing of whole app to guarantee reliability.
\end{itemize}

\subsection{Deployment}
\begin{itemize}
    \item \textbf{Vercel} --- The frontend is deployed on Vercel, leveraging its CDN for fast global delivery and integration with GitHub for continuous deployment.
\end{itemize}


\subsection{Industry Resources}
To justify the frontend technology choices, several authoritative industry resources were consulted:

\begin{itemize}
    \item \textbf{JavaScript (JS)}: Recognized as the most widely used programming language for web development, according to the Stack Overflow Developer Survey.  
    \url{https://survey.stackoverflow.co/2024/}

    \item \textbf{React}: Maintained by Meta and widely adopted in industry, React is considered a standard for building scalable, component-based user interfaces.  
    \url{https://react.dev/}

    \item \textbf{Vite}: A next-generation frontend tooling that provides fast Hot Module Replacement (HMR) and optimized builds, recommended by many open-source contributors and adopted across projects for its developer experience.  
    \url{https://vitejs.dev/}

    \item \textbf{CSS}: CSS is simple to use, and what was familiar to the group. No hassle, no need to learn syntax making the whole design experience seamless.
    \url{https://www.w3schools.com/css/}
\end{itemize}


\section{Backend Stack}
Overview:
This document provides a comprehensive guide to setting up and running the SDP Project backend API. The backend is built with Node.js, Express, Sequelize ORM, and PostgreSQL (NeonDB).

\textbf{TECHNOLOGY STACK FOR BACKEND}

\begin{itemize}
    \item \textbf{Runtime:} Node.js
    \item \textbf{Framework:} Express.js
    \item \textbf{Database:} PostgreSQL (NeonDB)
    \item \textbf{ORM:} Sequelize
    \item \textbf{Documentation:} Swagger
    \item \textbf{Authentication:} JWT (JSON Web Tokens)
    \item \textbf{Development Tool:} Nodemon
\end{itemize}

\textbf{PREREQUISITES} 
\begin{itemize}
    \item \textbf{Node.js:} v16 or higher
    \item \textbf{Package Manager:} npm
    \item \textbf{Database:} PostgreSQL (NeonDB recommended)
    \item \textbf{Version Control:} Git
\end{itemize}


\subsection{Database System}
\textbf{Option A: Local PostgreSQL}
\begin{itemize}
    \item \textbf{Install PostgreSQL locally}
    \item \textbf{Create a database}
    \item \textbf{Update \texttt{.env} with local credentials}
\end{itemize}

\textbf{Option B: NeonDB (Recommended)}
\begin{itemize}
    \item \textbf{Sign up at} \url{https://neon.tech}
    \item \textbf{Create a new project}
    \item \textbf{Copy the connection string}
    \item \textbf{Update \texttt{.env} with NeonDB credentials}
\end{itemize}


\subsection{Industry Resources}
\textbf{BACKEND References used:}

\begin{enumerate}
  \item \url{https://neon.com/docs/introduction} — official Neon docs. 
  \item \url{https://www.postgresql.org/docs/} — official Postgres manual. 
  \item \url{https://swagger.io/docs/} — OpenAPI / Swagger official site. 
  \item \url{https://en.wikipedia.org/wiki/Express.js} — details on Express framework. 
  \item \url{https://www.npmjs.com/package/nodemon} — nodemon npm page. 
  \item \url{https://www.youtube.com/watch?v=IlPPDIRjBKs} — YouTube video tutorial. 
  \item \url{https://www.youtube.com/watch?v=SccSCuHhOw0} — quick Express.js tutorial. 
  \item \url{https://www.youtube.com/watch?v=dhMlXoTD3mQ} — Swagger video. 
  \item \url{https://www.youtube.com/watch?v=f2EqECiTBL8} — comprehensive Node.js + Express course. 
  \item \url{https://youtu.be/9ZixI8Xy0T0?si=8YIfUflRKnW7CmtN} - postgres quick Overview
\end{enumerate}


\section{Testing Approach}
[Testing philosophy and coverage goals]

We aim to best the APIs using vitest because it fast and simple.
while for frontend e2e testing os more reliable. I covers most cases and most integration. 

\subsection{Unit Testing}
[Testing tools and methodology]
Backend - vitest

\subsection{End-to-End Testing}
[UI automation strategy]
Cypress - frontend testing Replaced by Playwright due to the multi window feature

\subsection{Industry Resources}
[Testing ROI studies or standards]
\url{https://testgrid.io/blog/cypress-testing/}
\url{https://www.testim.io/blog/cypress-vs-selenium/}
\url{https://dev.to/walmyrlimaesilv/10-reasons-why-you-should-use-cypress-for-web-testing-automation-jlb}
\url{https://www.lambdatest.com/blog/cypress-vs-playwright/}
\url{https://www.browserstack.com/guide/playwright-vs-cypress}
\url{https://www.checklyhq.com/learn/playwright/playwright-vs-cypress/}
\url{https://testgrid.io/blog/cypress-vs-playwright/}

\section{Development Tools}
[Team development workflow]

\subsection{Version Control}
[Git workflow and branching model]

\subsection{CI/CD Pipeline}
[Automation tools and stages]

\subsection{Hosting Platform}
\textbf{Backend hosted on Render}

\subsection{Industry Resources}
[Reliability engineering standards]
\url{https://www.conventionalcommits.org/en/v1.0.0/}
\url{https://blog.devgenius.io/make-a-meaningful-git-commit-message-with-semantic-commit-message-b39a79b13aa3}